% !TEX root = updates.tex
\begin{addblock}
\subsection{The best-response MIP for the Workshop 2 fleets}\label{sec:mipflow}
The program above is exact, but it does not scale to the Workshop 2 fleets. This subsection shows why and gives the
version we use. The rules of the game (\S\ref{sec:actions}--\S\ref{sec:payoffs}) do not change. Only the way the
player's decision problem is written down changes, and it stays exact.

\textbf{Why the Workshop 1 program breaks:} it lists every count vector a zone could end up with. If $A_{k,v}$ ships
of type $k$ could arrive at zone $v$ this period, the number of options in that zone is
\[
|\{b_{v,(n,p)}\}| \;=\; 2\prod_{k\in\K_i}\big(A_{k,v}+1\big).
\]
In Workshop 1, $|\K_\B|=1$ and $|\K_\Rd|=2$, so this was a few hundred. With the eight Blue types at full size and, for example,
$A_{k,v}=(12,13,16,12,2,4,9,6)$, it is
$2\cdot13\cdot14\cdot17\cdot13\cdot3\cdot5\cdot10\cdot7\approx8\times10^{7}$
options in one zone. The model cannot be built, whatever the license or the run time.

\textbf{Step 1, moves as flows:} free ships of the same type, at the same node, with the same set of allowed
destinations are interchangeable. Call such a set a group $g$, with type $k(g)$, location $u(g)$, allowed
destinations $\mathcal W_g$ and size $|g|$. Replace the per-ship binaries $y_{j,w}$ by integer flows
\[
f_{g,w}\in\Z_{\ge0}\ \ (w\in\mathcal W_g),\qquad \textstyle\sum_{w\in\mathcal W_g}f_{g,w}=|g| .
\]
These are the flows $f^t_{k,uv}$ of \S\ref{sec:actions}. Every assignment of ships gives a flow and every flow gives an
assignment, so this step changes nothing.

\textbf{Step 2, options by navy and coast guard strength level:} a zone's value depends on the player's ships there only
through two numbers, its navy strength and its coast guard strength. Control uses their sum. The assert hazard also
uses the navy share, because it blends the navy and coast guard values,
$h^i_{\As,v}=\theta_{i,v}h^i_{\As,\text{navy}}+(1-\theta_{i,v})h^i_{\As,\text{cg}}$ (\S\ref{sec:dynamics}).
Every weight $s\,w_k$, halved or not in rough seas, is a whole multiple of the unit $u=0.05\,s$ ($u=0.15$ for $s=3$), so
$u_k(w_r)=s\,w_k(w_r)/u$ is an integer. The navy and coast guard levels in zone $v$ are
\[
n_v=\sum_{k\ \text{navy}}u_k\Big(\bar x_{k,v}+\sum_{g:\,k(g)=k}f_{g,v}\Big),\qquad
c_v=\sum_{k\ \text{coast guard}}u_k\Big(\bar x_{k,v}+\sum_{g:\,k(g)=k}f_{g,v}\Big),
\]
so the own un-asserted strength is $R_v=u\,(n_v+c_v)$ and the navy share is $\theta_{i,v}=n_v/(n_v+c_v)$. The option
$b_{v,(n,c,p)}\in\{0,1\}$ means ``navy level $n$, coast guard level $c$ and posture $p$ in zone $v$''. Only levels that the
fixed ships and the arriving groups can actually produce are listed. The option values are the same formulas as before,
\[
a_v(n,c,p)=V^i_v\,\beta\,\sigma_i\phi_v(n,c,p)-\Lambda\rho_v(n,c,p),\qquad S_{i,v}=\alpha^{\1[p=\As]}\,u\,(n+c),
\]
with the hazard in $\rho_v$ using $\theta_{i,v}=n/(n+c)$. Both control and incident risk are therefore exactly those of the
Workshop 1 program. A patrol option does not depend on the split between $n$ and $c$, because the navy share enters only the
player's own assert hazard. So patrol options need only the total level, which keeps the program small.

\textbf{The program:}
\begin{align*}
\max_{f,\,b}\ \ &\textstyle\sum_{v\in\C}\sum_{(n,c,p)}a_v(n,c,p)\,b_{v,(n,c,p)}
-\gamma\sum_{g}\sum_{w\ne u(g)}f_{g,w}-1000\,\xi
&&\text{(control value $-$ incident cost $-$ fuel cost)}\\
\text{s.t.}\ \ &\textstyle\sum_{w\in\mathcal W_g}f_{g,w}=|g|\quad\forall g
&&\text{(every free ship gets one destination)}\\
&\textstyle\sum_{g}\sum_{w\ne u(g)}f_{g,w}\le F_i-\text{(forced moves)}
&&\text{(movement budget)}\\
&\textstyle\sum_{(n,c,p)}b_{v,(n,c,p)}=1\quad\forall v
&&\text{(one option per zone)}\\
&\textstyle\sum_{(n,c,p)}n\,b_{v,(n,c,p)}=n_v,\quad \sum_{(n,c,p)}c\,b_{v,(n,c,p)}=c_v\quad\forall v
&&\text{(the option matches the moves)}\\
&\textstyle\sum_{v\in\C}\sum_{(n,c,p)}\rho_v(n,c,p)\,b_{v,(n,c,p)}\le\varepsilon_i\,(1-\ell^t/L)+\xi
&&\text{(escalation-risk budget)}\\
&f\ \text{integer},\ b\ \text{binary},\ \xi\ge0.
\end{align*}
Compared with \S\ref{sec:mip}, $y$ becomes $f$, and the index $n$ (a count for each type) becomes $(n,c)$ (a navy and a
coast guard strength level). Everything else is unchanged. In the code, the options of one posture (and, for assert, one coast guard
level) are entered as a piecewise-linear function of the level with a breakpoint at every level that can occur, which is
the same program with fewer variables.

\textbf{Check:} at every solve we checked, the program's objective equals the value of its plan recomputed by the
game engine (\texttt{plan\_value}), and it is never below the value of the greedy plan for the same state.

\textbf{Run time and the game-ship scale:} at full size, even an earlier version of this program that approximated the
navy share took 13 seconds to 4 minutes per solve, and a one-year game needs 360 solves (two players, 180 periods).
With each game ship standing for three real ships (\S\ref{sec:hullscale}), a year of MIP-BR against MIP-BR takes 30 to 39
minutes on one core. The experiments (about 12{,}400 games) therefore run as independent jobs on MIT's Engaging cluster.
\end{addblock}
